\section{Frattini coordinates}\label{frattini-coordinates}

\textbf{Theorem UA.Phi (primitive saturation and coordinates).} Let
\(P\) be a finite \(p\)-group with an irredundant generating marking
\((r_i)_{i\in I}\), \(|I|=d\), and let \(F=\Phi(P)\). Then, for every
\(S\subseteq I\), \[
\Phi(P_S)=P_S\cap F,\qquad
P_SF=\pi^{-1}\operatorname{span}(\bar r_i:i\in S),\qquad
[P:P_SF]=p^{d-|S|}.
\tag{P1}
\] The saturated parabolics \(P_SF\) intersect according to intersection
of index sets. Each omission subgroup \(P_{I\setminus\{i\}}\) has
exactly one maximal overgroup, namely \(P_{I\setminus\{i\}}F\).

If \(B\geq F\) and \((s_i)_{i\in J}\) is an irredundant relative
generating family over \(B\), then \(B\unlhd P\), \(P/B\) has basis
\((s_iB)_{i\in J}\) over \(\mathbb F_p\), and \[
\langle B,s_i:i\in T\rangle/B=\operatorname{span}(s_iB:i\in T).
\tag{P2}
\] If \(p=2\) and all original \(r_i\) are involutions, then \(F=P'\)
and \(\Phi(P_S)=P_S'\) for every \(S\).

\emph{Proof.} Burnside's basis theorem makes the \(\bar r_i\) a basis of
\(P/F\). The image of \(P_S\) therefore has order \(p^{|S|}\), with
kernel \(P_S\cap F\). Since this image is elementary abelian,
\(\Phi(P_S)\leq P_S\cap F\). On the other hand \(P_S\) is generated by
\(|S|\) elements, so Burnside gives \([P_S:\Phi(P_S)]\leq p^{|S|}\). The
quotient by \(P_S\cap F\) already has that order. Equality follows,
proving the first formula. The remaining formulas and the intersection
assertion are inverse-image statements for coordinate subspaces.

Every maximal subgroup of a finite \(p\)-group contains \(F\) and
corresponds to a hyperplane of \(P/F\). A hyperplane containing a
specified coordinate hyperplane equals it, proving uniqueness of the
maximal overgroup of an omission subgroup.

For the relative assertion, every subgroup containing \(F\) is normal,
because \(P/F\) is abelian. The quotient \(P/B\) is elementary abelian.
The indicated images span it, and any nontrivial linear relation would
express one \(s_i\) using \(B\) and the other markers, contradicting
irredundancy. They form a basis, proving (P2). Finally, an abelian group
generated by involutions is elementary abelian of exponent two. Thus
\(P'\leq\Phi(P)\), while the elementary abelian quotient \(P/P'\) gives
the reverse inclusion. The same argument applies to each \(P_S\).
\(\square\)

These coordinates also describe the corresponding omission cosets: a
coset prescribes one coordinate value, and every simultaneous assignment
gives a chamber. The geometry is complete multipartite, with \(p\)
objects in each type, and the quotient translation group acts regularly
on chambers. For \(p>2\) it need not be the full type-preserving
automorphism group; the thin case \(p=2\) is the one covered by
UA.Geometry.
